\section{Background and Motivation}
\label{sec:background}

\subsection{5G Slice Management}

3GPP describes a network slice instance as an end-to-end logical network composed of slice subnets across domains such as the core network, RAN, and transport network~\cite{3gpp_ts28530}. Management systems commonly decompose this responsibility into a network slice management function (NSMF) and domain-specific network slice subnet management functions (NSSMFs)~\cite{3gpp_ts28533}. This decomposition is useful, but it leaves the operator with the problem of mapping service objectives into domain-specific parameters.

For the core network, free5GC implements a 5G core following the service-based architecture of 3GPP Release 15~\cite{free5gc,3gpp_ts23501}. Relevant slice and QoS controls include S-NSSAI, subscriber data, session management, QoS flow identifiers, and 5QI. For the RAN, OCUDU exposes practical controls for slice-aware scheduling through S-NSSAI-specific parameters, including PRB policy ratios and priority~\cite{ocudu,ocudu_oran_gnb_doc}. The challenge is not the absence of knobs, but the lack of a unified layer that derives them from an operator intent.

\subsection{Kubernetes, Operators, and Nephio}

Kubernetes offers a declarative API and a reconciliation model that are well matched to intent management. Custom Resource Definitions (CRDs) extend the API with domain-specific objects~\cite{kubernetes_custom_resources}. Operators then reconcile desired state into actual runtime state~\cite{kubernetes_operator_pattern,kubernetes_controller}. This makes Kubernetes a natural substrate for expressing both high-level intent and low-level network-function configuration.

Nephio extends this model with KRM packages, Porch package revisions, and GitOps synchronization~\cite{nephio,google_config_sync,cncf_gitops_2025}. This provides version control, review, and rollback for network configuration. Nephio provides a free5GC Operator that reconciles \texttt{NFDeployment} resources and manages the lifecycle of core network functions. The RAN path uses the OCUDU Operator to perform the corresponding domain-specific translation and reconciliation for OCUDU resources.

\subsection{Gap}

Prior work on intent-based and autonomous network management proposes architectures for translating service objectives into network behavior~\cite{chirivella2023e2e,martini2023intent,tran2024ina,antonakoglou2025camino,nahum2025intent}. Recent work also studies GitOps performance~\cite{ghosh2025performance}. The remaining gap is an end-to-end prototype connecting intent submission, cross-domain translation, Nephio package workflows, RAN and core NSSMF reconciliation, free5GC QoS configuration, and experimental validation.

\begin{table}[t]
  \caption{Positioning of this work relative to open-source slice-management building blocks.}
  \label{tab:positioning}
  \begin{tabular}{p{0.20\linewidth}p{0.30\linewidth}p{0.34\linewidth}}
    \toprule
    Component & Primary role & Limitation addressed here \\
    \midrule
    free5GC & 5G core and QoS configuration & No end-to-end intent interface \\
    OCUDU & RAN and slice scheduler parameters & Requires low-level scheduler configuration \\
    Kubernetes & Declarative API and reconciliation & Needs telecom-specific translation logic \\
    Nephio & Package lifecycle and free5GC Operator & Needs intent-derived RAN mutation \\
    \bottomrule
  \end{tabular}
\end{table}

Table~\ref{tab:positioning} highlights the design premise. Each component solves part of the deployment problem, but none provides a complete path from service intent to realized RAN and core configuration. The E2E orchestrator connects their control surfaces; the OCUDU and free5GC Operators then serve as the RAN and core domain reconcilers, respectively.
